\section{Stronger algorithms for small caps}
\label{sec:small-caps}

Consider instances with a distinguished hypothesis $p$ and orient
each supplied query so that its \emph{cap} is the answer side missing
$p$. Let $P$ be the other $N$ hypotheses. These assumptions concern
the supplied queries themselves. At a recursive flat state, the
common external witness serves as an analysis aid, while $p$ remains
the distinguished hypothesis.
The restricted algorithms below clarify which small-cap constructions
can support approximation hardness.

\subsection{Exact maximum depth with bounded caps}

\begin{theorem}\label{thm:bounded-cap-height}
If every cap has size at most a fixed integer $b$, maximum depth is
exactly solvable in polynomial time. A running-time bound is
$m^{b-1}2^{O(b^2)}\operatorname{poly}(n,m,L)$, where $L$ is the
input bit length.
\end{theorem}

\begin{proofstrategy}
\begin{enumerate}
\item A minimum cap cover gives a lower bound $K$ from the
witness spine and a tree of height at most $K+b-1$.
\item Only the last $b-1$ witness queries need to be guessed.
Their caps contain at most $b(b-1)$ points, so guess their residual
state and find the earlier cover exactly.
\item Exact local subset tables verify every suffix exit. A feasible
optimal tree supplies one of these guesses, proving exactness.
\end{enumerate}
\end{proofstrategy}

\subsubsection{Supporting propositions and the algorithm}

Let $K$ be the minimum number of supplied caps covering $P$. For a
target height $D\in[K,K+b-1]$, let $d_0=\max\{0,D-b+1\}$.
A suffix guess consists of at most $D-d_0$ ordered witness queries
and the set $R\subseteq P$ alive immediately before them.

The witness path must eliminate every other hypothesis, while each
positive exit contains only a few hypotheses. Comparing these two
requirements confines the optimum to a short interval.

\begin{proposition}[The optimum lies in a short height interval]
\label{prop:cap-height-range}
A minimum supplied-halfplane cover is computable in polynomial time,
and the optimal height lies in $[K,K+b-1]$.
\end{proposition}

Early exits have enough remaining depth for any proper local code.
The only delicate exits occur near the end of the witness path;
preserving their guessed residual set makes the earlier cover independent.

\begin{proposition}[Each suffix guess is exactly checkable]
\label{prop:cap-height-suffix}
Guess $R$ inside the suffix caps' union. A prefix cover of
$P\setminus R$ using caps disjoint from $R$, followed by the suffix,
gives height at most $D$ if its cover size is at most $d_0$ and every
suffix first-hit group $A_j$ satisfies
$j+h(A_j)\le D-d_0$, where $j$ is its first-hit position within the
suffix and $h$ is exact local identification height.
These conditions are decidable in polynomial time for fixed $b$.
\end{proposition}

Exact checking is useful only if an optimum appears among the guesses.
Cutting an optimal witness path at the prescribed depth provides
precisely the small suffix and residual set being enumerated.

\begin{proposition}[Enumeration is complete and polynomial]
\label{prop:cap-height-complete}
Every height-$D$ tree provides an accepted suffix guess. There are
$O(m^{b-1})$ ordered suffixes and $2^{O(b^2)}$ candidate residual
sets per suffix. For $b=1$, a minimum cover itself is exact.
\end{proposition}

\begin{algorithm}[H]
\caption{Exact maximum depth with a distinguished miss hypothesis and bounded caps}
\label{alg:bounded-caps}
\KwIn{Distinguished hypothesis $p$; other hypotheses $P$; supplied caps of size at most fixed $b$.}
\KwOut{An exact minimum-height supplied-query tree.}
\If{$P=\varnothing$}{\KwRet the zero-query tree\;}
Compute the minimum supplied-halfplane cover size $K$ of $P$\;
\For{$D=K,\ldots,K+b-1$}{
  $d_0\gets\max\{0,D-b+1\}$\;
  \ForEach{ordered supplied-query suffix of length $q\le D-d_0$}{
    \ForEach{$R$ contained in the union of its caps}{
      \If{the suffix does not cover $R$}{\Continue\;}
      Find a minimum cover $Q$ of $P\setminus R$ using only caps disjoint from $R$\;
      \If{$|Q|>d_0$ or the cover is infeasible}{\Continue\;}
      Compute the suffix's first-hit groups $A_j$ on $R\cup\{p\}$\;
      \If{some $j+c_W(\ExactTree(A_j,\mathrm W))>D-d_0$}{\Continue\;}
      Execute $Q$ followed by the suffix on the all-miss spine\;
      Attach exact local supplied-query trees to all positive first-hit groups\;
      \KwRet this tree\;
    }
  }
}
\end{algorithm}

The theorem's height interval means the loop needs only the listed
values of $D$. For each value, the algorithm searches suffixes and
residual sets; since a suffix has at most $b-1$ queries and each cap
contains at most $b$ points, its union contains at most $b(b-1)$
points. An empty suffix therefore permits only $R=\varnothing$. The
prefix step is a minimum supplied-halfplane cover, computed by the
polynomial algorithm cited in Theorem~\ref{thm:bounded-cap-height};
when the instance has the special hull-interval representation, its
dedicated interval oracle performs the same step.


\begin{proof}[Proof of Theorem~\ref{thm:bounded-cap-height}]
For $P=\varnothing$, return the zero-query tree. Otherwise,
Proposition~\ref{prop:cap-height-range} limits the target heights.
Proposition~\ref{prop:cap-height-suffix} proves that every accepted
guess constructs a legal tree, and
Proposition~\ref{prop:cap-height-complete} ensures acceptance at the
optimum. Returning the first accepted height is exact; the stated
enumeration and polynomial cover and local optimizations give the
running time.
\end{proof}

\subsubsection{Technical proofs of the height propositions}

\begin{proof}[Proof of Proposition~\ref{prop:cap-height-range}]
Let $K$ be the minimum number of supplied caps covering $P$.
Minimum planar halfplane cover is exactly solvable in polynomial
time~\cite{harpeledlee2012}. Separability ensures that the cover
exists. The path identifying $p$ in any tree must cover $P$, so
its height is at least $K$. Conversely, query a minimum cover along
the all-miss spine and identify each positive hit group locally.
A hit group has at most $b$ members and therefore admits a proper
identification tree of height at most $b-1$. This gives height at
most $K+b-1$. Only these $b$ target heights need to be tested.

\end{proof}

\begin{proof}[Proof of Proposition~\ref{prop:cap-height-suffix}]
For a target $D\in[K,K+b-1]$, put $d_0=\max\{0,D-b+1\}$.
Guess an ordered suffix of $q\le D-d_0\le b-1$ supplied queries
on the witness spine. Guess also the set $R\subseteq P$ remaining
just before that suffix. Require that the suffix caps cover $R$;
thus $R$ is a subset of their union and has size at most $b(b-1)$.
For each guess, compute a minimum cap cover of $P\setminus R$
using only supplied caps disjoint from $R$, and require size at
most $d_0$. An empty suffix permits $R=\varnothing$.

The suffix must fit on both kinds of path: a positive exit needs
its local code, and the all-miss route must also terminate within
the budget. Checking both prevents a hidden witness-path overrun.

\begin{claim}[Local suffix tests include the witness leaf]
The exact checks verify every positive exit and the full miss route.
\end{claim}
\begin{proof}
Execute the guessed suffix on $R\cup\{p\}$. For each nonempty
first-hit group $A_j$ at position $j$, compute its exact local
identification height $h(A_j)$ by subset Bellman DP on at most
$b$ hypotheses. Require $j+h(A_j)\le D-d_0$ and $q\le D-d_0$.
These tests verify the entire suffix tree, including the witness
leaf. The prefix cover takes at most $d_0$ queries, preserves $R$
on its all-miss route, and its positive exits cost at most $b-1$
additional queries. Thus its exits have height at most
$d_0+b-1=D$ when $d_0>0$; if $d_0=0$, the prefix is empty.
The equality here uses the positive branch of
$d_0=\max\{0,D-b+1\}$; no such equality is needed when the prefix
has length zero.
The suffix starts no later than depth $d_0$, so an accepted guess
constructs a tree of height at most $D$.

\end{proof}
\end{proof}

\begin{proof}[Proof of Proposition~\ref{prop:cap-height-complete}]
For completeness, take an identifying tree of height at most $D$. If its witness
path ends before depth $d_0$, use an empty suffix and $R=\varnothing$;
its witness queries cover $P$ within $d_0$. Otherwise take the
remaining witness queries after depth $d_0$ as the suffix, and the
actual nonwitness state there as $R$. The prefix caps are disjoint
from $R$ and cover $P\setminus R$. The suffix has at most $b-1$
queries and covers $R$; its positive subtrees satisfy the exact
local-height tests. This guess is accepted. There are
$O(m^{b-1})$ suffix descriptions and $2^{O(b^2)}$ candidate subsets
per suffix, and all remaining computations are polynomial for
fixed $b$. For $b=1$, a minimum cover itself is exact.
\end{proof}

\subsection{An EPTAS for weighted singleton and pair caps}

Assume every cap has size one or two, all singleton caps for $P$
are supplied, and the graph $G$ of two-point caps is a disjoint
union of paths and cycles. The latter condition follows for planar
instances with all singleton caps: every nonwitness point is a
convex-hull vertex, and a two-point cap joins consecutive vertices.
Let the positive raw weights be $w_0$ for $p$ and $w_v$ for $v\in P$,
and put $\Omega=w_0+\sum_v w_v$.

For a matching $M$, replace each matched pair by a group of mass
$w_u+w_v$ and leave unmatched vertices as singleton groups. Let
$J(M)$ be the group multiset, $k=N-|M|$, and
$P(M)=\sum_{uv\in M}(w_u+w_v)$. With group masses
$a_1\ge\cdots\ge a_k$, write $C(J)=\sum_j j a_j$. The matching
cost is
\begin{equation}\label{eq:eptas-matching-invoice}
 G(M)=w_0k+C(J(M))+P(M).
\end{equation}

\begin{theorem}\label{thm:pair-eptas}
For this weighted singleton/pair path-and-cycle class, expected
depth has a deterministic $(1+\varepsilon)$ approximation in time
$2^{O(\varepsilon^{-3})}\operatorname{poly}(n,L,\log(1/\varepsilon))$
for $0<\varepsilon\le1$. The polynomial exponent is independent
of $\varepsilon$; this is an EPTAS.
\end{theorem}

\begin{proofstrategy}
\begin{enumerate}
\item Delete one residue class of path and cycle edges. Averaging
finds a near-optimal matching on bounded-size components.
\item Enumerate components containing the heaviest vertices. Their
complement has little mass, so only light components need rounding.
\item Express the exact matching cost as a convex quadratic
program on component configurations. Its fractional optimum
lower-bounds the best integral invoice; the computed feasible solution
has only an additive numerical optimization error.
\item Independent configuration rounding adds only a controlled
variance term. Conditional expectations make it deterministic.
\end{enumerate}
\end{proofstrategy}

\subsubsection{Supporting propositions and the algorithm}

Work with raw weights and put $F^*=\min_MG(M)=\Omega\OPT_E$.
The graph, weights, group counts, and count functions below describe
the instance. The tuning controls are integers $K\ge1$, $H\ge0$ and
an additive accuracy $0<\eta\le1$, chosen only in the final calibration.
Put $B=2K$, the resulting component-size bound. Take the $H$ heaviest
nonwitness vertices, or all of them if fewer than $H$ exist, and let
$\Omega_{\rm tail}$ be the mass of the others.

Before approximating matchings, we need to justify that they describe
the full adaptive problem. The witness spine supplies this link;
the supplied singleton queries realize every resulting matching.

\begin{proposition}[The matching cost formula is exact]
\label{prop:pair-invoice-exact}
Minimizing \eqref{eq:eptas-matching-invoice} is exactly the minimum
raw weighted depth, including arbitrary adaptive trees. Every matching
is executable with supplied singleton and pair queries.
\end{proposition}

Two reductions make the remaining optimization manageable. Edge
deletion bounds component size at a small relative cost, while
fixing heavy components leaves little mass exposed to rounding.

\begin{proposition}[Residue deletion and heavy enumeration]
\label{prop:pair-deletion-heavy}
Deleting edge numbers in one residue modulo $K$ leaves components of
at most $B$ vertices. Some residue has a surviving matching of cost
at most $(1+1/K)F^*$. Moreover,
$\Omega_{\rm tail}\le2F^*/(H+1)$, and all matching choices on components
meeting a heavy vertex can be enumerated in at most $2^{BH}$ trials.
\end{proposition}

For a fixed heavy trial, enumerate the matching configurations $c$ of
each remaining, light component $i$. Let $k_{ic}$ be its group count,
$P_{ic}$ its matched-pair mass, and $z_{ic}(t)$ its number of groups
of mass at least $t$. Let $Z_f(t),k_f,P_f$ be the corresponding totals
of the fixed components. Fractional configuration weights satisfy
$x_{ic}\ge0$ and $\sum_cx_{ic}=1$. Set
$\mu(t)=Z_f(t)+\sum_{ic}x_{ic}z_{ic}(t)$ and minimize
\begin{equation}\label{eq:pair-convex-relaxation}
 Q(x)=w_0\left(k_f+\sum_{ic}x_{ic}k_{ic}\right)
       +P_f+\sum_{ic}x_{ic}P_{ic}
       +\frac12\int_0^\infty(\mu(t)^2+\mu(t))\,dt.
\end{equation}

Group order couples choices made in different components. Counting
groups above each mass threshold captures this coupling in one
quadratic objective, without guessing the final order.

\begin{proposition}[An exact convex quadratic model]
\label{prop:pair-convex-model}
The displayed objective equals the matching cost on integral
configurations. Its minimum over the product of configuration simplices
lower-bounds the best integral invoice. It is a rational convex quadratic program
with $N2^{O(B)}$ variables and $O(N)$ mass breakpoints. A feasible
rational solution within raw additive error $\eta\Omega$ is computable
with polynomial bit complexity and time polynomial in the model size,
$L$, and $\log(1/\eta)$.
\end{proposition}

Rounding combines the component choices into one matching. Only the
light components fluctuate, and their mass bounds the extra cost of
turning the relaxation into a tree.

\begin{proposition}[Deterministic rounding costs only light mass]
\label{prop:pair-rounding}
From a feasible rational $x$, conditional expectations produce a matching of
cost at most $Q(x)+(B/2)\Omega_{\rm tail}$ in polynomial time in the
model size and bit length. If every trial uses a feasible rational $x$
within $\eta\Omega$ of its fractional optimum, the cheapest tree over all
residues and heavy trials has cost at most
\[
 \left(1+\frac1K+\eta+\frac B{H+1}\right)F^*.
\]
\end{proposition}

\begin{algorithm}[H]
\caption{Expected-depth EPTAS for weighted singleton/pair paths and cycles}
\label{alg:pair-eptas}
\KwIn{Distinguished miss hypothesis, supplied singleton tests and pair graph $G$;
positive raw weights; $0<\varepsilon\le1$.}
\KwOut{A deterministic $(1+\varepsilon)$ approximation.}
Use the controls $K,H,\eta$ from
Section~\ref{sec:pair-eptas-calibration}; set $B\gets2K$\;
\If{$n=1$}{\KwRet the zero-query tree\;}
Find the $H$ heaviest nonwitness vertices (or all of them if fewer exist)\;
\For{$r=0,\ldots,K-1$}{
  Delete edges numbered $r$ modulo $K$ on each path or cycle\;
  Enumerate each remaining component's matching configurations\;
  \ForEach{joint configuration on components meeting the heavy vertices}{
    Fix these heavy configurations\;
    Solve the light-component convex program \eqref{eq:pair-convex-relaxation}
    to raw additive accuracy $\eta\Omega$\;
    \ForEach{unfixed light component}{
      Fix a configuration minimizing the conditional expected invoice,
      using rational first and second count moments\;
    }
    Form the matched-pair and singleton groups; sort them by decreasing mass\;
    Query the supplied group tests on the miss spine,
    using one supplied singleton test to identify each pair hit\;
    Retain the resulting tree and its exact expected cost\;
  }
}
\KwRet the cheapest retained tree\;
\end{algorithm}

Deleting one edge residue bounds the component size, and enumerating
the heavy components leaves a light-instance relaxation. The
conditional-expectation step then fixes each remaining component while
preserving the invoice bound, so every retained candidate is an actual
matching and hence yields a legal tree. Propositions
\ref{prop:pair-deletion-heavy}--\ref{prop:pair-rounding} prove the
component bound, heavy enumeration, convex relaxation, and rounding
guarantee. Proposition~\ref{prop:pair-invoice-exact} then converts the
rounded matching into its exact tree cost. The fractional solution
guides the rounding, while the returned tree comes from an actual
matching.


\begin{proof}[Proof of Theorem~\ref{thm:pair-eptas}]
For $n=1$, return the exact zero-query tree. For $n\ge2$, choose
$K,H,\eta$ by Section~\ref{sec:pair-eptas-calibration}.
Proposition~\ref{prop:pair-deletion-heavy} provides a good residue and
an enumerated heavy part of its near-optimal matching.
Propositions~\ref{prop:pair-convex-model} and~\ref{prop:pair-rounding}
return cost at most $(1+\varepsilon)F^*$.
Proposition~\ref{prop:pair-invoice-exact} turns the matching into a
legal supplied-query tree. Divide by $\Omega$ to obtain the stated
expected-depth ratio. The calibration gives $BH=O(\varepsilon^{-3})$;
the other stages have polynomial exponent independent of $\varepsilon$.
\end{proof}

\subsubsection{Technical proofs of the EPTAS propositions}

\begin{proof}[Proof of Proposition~\ref{prop:pair-invoice-exact}]
Follow the witness in an arbitrary proper tree. Its first-hit groups
are disjoint singletons or supplied pairs, so the pair groups form
a matching $M$ and the number of groups is $k=N-|M|$. A supplied pair
query first hitting only one vertex can be replaced here
by that vertex's supplied singleton test. The witness contributes
$w_0k$, group positions contribute $\sum_j j a_j$, and each pair
exit needs exactly one further supplied singleton query, contributing
its full pair mass. This is the displayed invoice before sorting.

For adjacent group masses $a<b$, putting $b$ first reduces their
contribution by
$[ja+(j+1)b]-[jb+(j+1)a]=b-a>0$. Thus nonincreasing order minimizes
the position term, giving $C(J(M))$. Conversely, query any matching's
disjoint pairs and unmatched singletons in that order, and attach
one supplied singleton query at each pair exit. This attains
$w_0k+C(J(M))+P(M)$, proving exactness for arbitrary trees.

\end{proof}

\begin{proof}[Proof of Proposition~\ref{prop:pair-deletion-heavy}]
We first control the cost of simplifying the graph, then control
the mass left for rounding. For graph simplification, compare each
residue with the same optimal matching by splitting its deleted pairs.

\begin{claim}[One residue loses at most a $1/K$ fraction]
\label{clm:pair-residue-loss}
Splitting the removed matched edges leaves a legal near-optimal matching.
\end{claim}
\begin{proof}
Number edges consecutively on each path or cycle. For each common
residue $r\in\{0,\ldots,K-1\}$, delete edges numbered $r$ modulo
$K$. A path component between consecutive deleted edges has at most
$K$ vertices, as do its end components. On a cycle, all components
except possibly the one crossing the numbering cut have the same
bound. The crossing component joins two end portions, each with fewer
than $K$ edges, and has at most $2K=B$ vertices. If a residue deletes
no edge of a cycle, that cycle has fewer than $K$ edges, so is already
smaller than $B$.

Take an optimal matching, with sorted group masses
$a_1,\ldots,a_k$. For a matched pair of weights $a\ge b$ at
position $j$, put $R_j=w_0+\sum_{l>j}a_l$. Its original weighted
depth contribution is $(j+1)(a+b)$, including the local pair query.
After replacing it by two consecutive singleton queries the contribution
is $ja+(j+1)b$, a change of $-a$. Every later group and the
witness are delayed by one query, a change of $R_j$. Thus the
net change is exactly $R_j-a$.

For simultaneous splits, each earlier split delays all points of a
later group by one regardless of whether that group is also split.
The local change $-a$ is likewise independent of the group's shifted
position. The net changes therefore add. Sorting the new singleton
and pair groups can only improve the temporary schedule. Since the
$R_j$ are positive,
\[
 \sum_{j\text{ matched}}R_j
 \le\sum_{j=1}^k R_j
 =kw_0+\sum_{l=1}^k(l-1)a_l
 \le kw_0+C(J)\le F^*.
\]
The first inequality adds the remaining nonnegative $R_j$ terms.
For the equality, group $l$ contributes to exactly $l-1$ earlier
tails. Replacing $l-1$ by $l$ gives the next inequality, and the
last discards the nonnegative pair-continuation term from the exact
invoice of Proposition~\ref{prop:pair-invoice-exact}.
Each matched edge belongs to exactly one residue trial. Averaging
the temporary schedules over the $K$ residues gives cost at most
\[
 F^*+\frac1K\sum_{j\text{ matched}}(R_j-a)
 \le F^*+\frac{F^*}{K}.
\]
This inequality first discards each nonpositive local saving $-a$
and then applies the preceding bound on the summed tails.
At least one trial meets this average bound, and sorting its groups
produces a surviving matching of no larger cost.

\end{proof}
A component can be small in vertex count yet carry a heavy vertex,
raising the rounding cost. Enumerating the heavy part removes this
risk, while the rank lower bound controls the remaining mass.

\begin{claim}[Heavy components leave a low-mass rounding problem]
\label{clm:pair-heavy-tail}
The rank lower bound controls all unfixed vertex mass.
\end{claim}
\begin{proof}
Order individual nonwitness weights as $w_{(1)}\ge\cdots\ge w_{(N)}$.
A first-hit group contains at most two vertices, so at most $2j$
vertices can leave by spine position $j$. Consequently the sorted
list of first-hit positions is componentwise at least
$\lceil1/2\rceil,\ldots,\lceil N/2\rceil$.
Pairing larger weights with smaller positions minimizes their weighted
sum by adjacent exchanges. Ignoring local continuation and witness
costs therefore gives
\[
 F^*\ge\sum_{i=1}^N\lceil i/2\rceil w_{(i)}.
\]
For every $i>H$, $\lceil i/2\rceil\ge(H+1)/2$, hence
\[
 F^*\ge\sum_{i>H}\lceil i/2\rceil w_{(i)}
 \ge\frac{H+1}{2}\sum_{i>H}w_{(i)}
 =\frac{H+1}{2}\Omega_{\rm tail}.
\]
The first inequality drops the nonnegative contributions of the
heaviest $H$ vertices, and the second applies the rank lower bound
uniformly to every remaining index. Thus the tail mass is bounded
without assuming any lower bound on an individual weight.
For each residue, enumerate all matching configurations in components
containing one of the $H$ heaviest vertices. There are at most $H$
such components. A path or cycle on at most $B$ vertices has at most
$B$ edges, so has at most $2^B$ edge subsets and hence at most that
many matching configurations. The joint number of choices is at most
$(2^B)^H=2^{BH}$. This enumeration
is independent of $n$ in its exponent. The remaining components'
total vertex mass is at most $\Omega_{\rm tail}$. If $H\ge N$,
all components are fixed and there is no rounding stage.

\end{proof}
\end{proof}

\begin{proof}[Proof of Proposition~\ref{prop:pair-convex-model}]
For sorted group masses $a_1\ge\cdots\ge a_k$,
the groups above a threshold occupy ranks $1$ through $Z(t)$.
Finite summation and integration therefore give
\[
 C(J)=\sum_{j=1}^k j a_j
     =\int_0^\infty\sum_{j:a_j\ge t}j\,dt
     =\int_0^\infty\frac{Z(t)(Z(t)+1)}2\,dt.
\]
For the first integral, integrate the indicator of $t\le a_j$
over $t\ge0$ and sum over $j$. The second integral then sums the
consecutive ranks $1,\ldots,Z(t)$. Consequently the formula prices
all cross-component ordering effects exactly.
On integral configurations, $\mu(t)=Z(t)$ and the linear
terms are exactly the witness and pair-continuation terms of the
invoice. Allowing fractional configurations enlarges the feasible
set, so its minimum cannot exceed the integral minimum.
The only breakpoints are singleton and available-pair masses, at most
$N+|E(G)|=O(N)$ rational values. Include zero among the ordered
breakpoints $\tau_j$. On each interval between consecutive breakpoints,
$\mu(t)$ is a fixed affine function $\mu_j(x)$, so the integral is
\[
 \frac12\sum_j(\tau_{j+1}-\tau_j)
                 \bigl(\mu_j(x)^2+\mu_j(x)\bigr).
\]
The coefficients are rational and the interval lengths are nonnegative;
thus its quadratic part is a nonnegative sum of squares of affine
functions and is convex. Each component has at most $2^B$ matching
configurations and there are at most $N$ components, giving
$N2^{O(B)}$ variables over a product of simplices. After normalizing by $\Omega$, solve it to additive
accuracy $\eta$ using polynomial-time convex quadratic
optimization~\cite{nestnem1994}. Equivalently the raw error is $\eta\Omega$. The simplex
constraints have rational coefficients and a feasible integral point
obtained by choosing one configuration per component.
For any feasible fractional point, $0\le\mu(t)\le N$ and
$\int\mu(t)\,dt=\sum_vw_v$, since every configuration partitions
its component's vertices. Its witness term is at most $Nw_0$,
its pair-continuation term is at most $\sum_vw_v$, and
\[
 \frac12\int(\mu(t)^2+\mu(t))\,dt
 \le\frac{N+1}{2}\int\mu(t)\,dt
 =\frac{N+1}{2}\sum_vw_v.
\]
Here $\mu(t)^2\le N\mu(t)$ follows pointwise from
$0\le\mu(t)\le N$, and the equality uses the invariant total group
mass across configurations.
Thus $0\le Q(x)\le(N+3)\Omega$, giving a normalized objective
bound $O(n)$. Together with bounded simplices and rational
coefficients of polynomial bit length, the cited convex-QP algorithm
provides the stated feasible rational approximation and bit complexity.

\end{proof}

\begin{proof}[Proof of Proposition~\ref{prop:pair-rounding}]
We separate the rounding analysis from its deterministic implementation.
The fractional objective records the mean counts; independent rounding
changes only their quadratic contribution, which can be charged to
light vertex mass.

\begin{claim}[Only the quadratic count term incurs variance]
\label{clm:pair-rounding-variance}
The expected rounded cost is at most $Q(x)+(B/2)\Omega_{\rm tail}$.
\end{claim}
\begin{proof}
Independently round each light component according to its fractional
configuration distribution, and write
$Z(t)=Z_f(t)+\sum_iX_i(t)$. Then $\mathbb EZ(t)=\mu(t)$.
For each component $0\le X_i(t)\le B$, so
\[
 \operatorname{Var}X_i(t)
 =\mathbb E[X_i(t)^2]-(\mathbb EX_i(t))^2
 \le\mathbb E[X_i(t)^2]\le B\mathbb EX_i(t).
\]
The first inequality removes a nonnegative square; the second uses
$X_i(t)^2\le BX_i(t)$ on $[0,B]$. Hence the variance bound depends
on component size and expected mass, regardless of how many
configurations the component has.
Independence gives
$\mathbb E[Z(t)^2]=\mu(t)^2+\sum_i\operatorname{Var}X_i(t)$.
The witness and pair-continuation terms are linear, so their expectations
are exactly their terms in $Q(x)$. Applying the rank integral from
Proposition~\ref{prop:pair-convex-model} and then the variance bound yields
\begin{align*}
 \mathbb EG(M)
 &=Q(x)+\frac12\int_0^\infty
                 \sum_i\operatorname{Var}X_i(t)\,dt\\
 &\le Q(x)+\frac B2\sum_i
                  \int_0^\infty\mathbb EX_i(t)\,dt.
\end{align*}
For every configuration, the integral of its count function is the
sum of its group masses. These groups partition the component's
vertices, so this sum is its fixed vertex mass, independent of the
configuration. The displayed sum of integrals is therefore the total
light-component mass, at most $\Omega_{\rm tail}$. Hence
\[
 \mathbb EG(M)\le Q(x)+\frac B2\Omega_{\rm tail}
                 \byref{Prop.}{prop:pair-deletion-heavy}{\le} Q(x)+\frac B{H+1}F^*,
\]
using the tail bound proved above.
\end{proof}
An expected-cost bound still needs an individual matching that attains
it. Fixing one component at a time preserves that bound and keeps
all intermediate computations rational and explicit.

\begin{claim}[Exact conditional expectations give a supplied tree]
One component at a time can be fixed without increasing the expected invoice.
\end{claim}
\begin{proof}
Suppose some components have already been fixed. For the next
component, the current conditional expected invoice is the weighted
average of the conditional expected invoices obtained by fixing each
configuration, with weights $x_{ic}$ summing to one. At least one
configuration of positive probability therefore has value no larger
than the current expectation. Choose a minimum such value and repeat.

Compute each value exactly from the fixed counts and the remaining
independent first and second count moments at the same $O(n)$
breakpoints. The identity for $\mathbb E[Z(t)^2]$ remains valid with
the fixed counts incorporated into $Z_f(t)$. All interval lengths,
probabilities, moments, and linear terms are rational, with polynomial
bit length; the number of configurations and breakpoints is the model
size. At the end no randomness remains, so the actual matching cost is
at most the initial expectation. Sorting its final groups yields
a supplied-query tree by Proposition~\ref{prop:pair-invoice-exact}.

\end{proof}
For the good residue, enumeration contains the heavy part of an optimum
surviving matching. The fractional optimum is no greater than that
matching's invoice. In a nontrivial instance every hypothesis has depth at least one,
so $F^*\ge\Omega$. For the good residue and heavy trial, the computed
feasible point thus satisfies
\[
 Q(x)\byref{Prop.}{prop:pair-convex-model}{\le} Q^*+\eta\Omega
       \byref{Prop.}{prop:pair-deletion-heavy}{\le}(1+1/K)F^*+\eta F^*.
\]
Claim~\ref{clm:pair-residue-loss} bounds the surviving comparison
invoice, while $F^*\ge\Omega$ converts the numerical error to a
relative one. Adding Claim~\ref{clm:pair-rounding-variance}, with the
tail estimate of Claim~\ref{clm:pair-heavy-tail}, gives cost at most
\[
 (1+1/K)F^*+\eta F^*+\frac B{H+1}F^*
 =\left(1+\frac1K+\eta+\frac B{H+1}\right)F^*.
\]
The algorithm takes the cheapest trial, so it meets this same bound.
\end{proof}

\subsubsection{Calibration of the EPTAS controls}
\label{sec:pair-eptas-calibration}

\begin{parameterchoice}
For $0<\varepsilon\le1$, choose
\[
 K=\lceil3/\varepsilon\rceil,\qquad B=2K,\qquad
 H=\lceil3B/\varepsilon\rceil,\qquad \eta=\varepsilon/3.
\]
Then $1/K+\eta+B/(H+1)\le\varepsilon$.
Also $B=O(\varepsilon^{-1})$ and $H=O(\varepsilon^{-2})$, so heavy
enumeration costs $2^{O(\varepsilon^{-3})}$. Configuration enumeration,
convex optimization, and conditional expectations keep a polynomial
exponent independent of $\varepsilon$.
\end{parameterchoice}

The ceiling choices give $K\ge3/\varepsilon$ and
$H+1\ge3B/\varepsilon$, hence
$1/K\le\varepsilon/3$ and $B/(H+1)\le\varepsilon/3$.
Together with $\eta=\varepsilon/3$ this proves the calibrated ratio.
For $0<\varepsilon\le1$, $K=O(\varepsilon^{-1})$, so
$B=O(\varepsilon^{-1})$ and
$H=O(B/\varepsilon)=O(\varepsilon^{-2})$; consequently
$BH=O(\varepsilon^{-3})$ as used in the enumeration bound.

All-singleton availability is essential both for the invoice's converse
and for splitting removed pairs. The EPTAS includes the weighted
pair-path hardness instances below, and its exponential dependence on
$1/\varepsilon$ is compatible with strong NP-completeness. This
subclass therefore cannot establish APX-hardness unless
$\mathrm P=\mathrm{NP}$.
